\documentclass[11pt]{article}

% ---------- Packages ----------
\usepackage[a4paper,margin=1in]{geometry}
\usepackage{amsmath,amssymb,amsthm,mathtools}
\usepackage{booktabs}
\usepackage{hyperref}
\usepackage{enumitem}
\usepackage[numbers]{natbib}

% ---------- Theorem environments ----------
\newtheorem{theorem}{Theorem}
\newtheorem{lemma}[theorem]{Lemma}
\newtheorem{corollary}[theorem]{Corollary}
\newtheorem{proposition}[theorem]{Proposition}

\theoremstyle{definition}
\newtheorem{definition}[theorem]{Definition}

\theoremstyle{remark}
\newtheorem*{remark}{Remark}
\newtheorem*{status}{Status}

% ---------- Commands ----------
\newcommand{\Lmu}{\mathcal{L}_{\mathrm{BA},\mu}}
\newcommand{\Q}{\mathbb{Q}}
\newcommand{\N}{\mathbb{N}}
\newcommand{\St}{\mathrm{St}}
\newcommand{\Prin}{\mathrm{Prin}}
\newcommand{\Supp}{\mathrm{Supp}_\sigma}
\newcommand{\pfa}{\mathrm{pfa}}
\newcommand{\fc}{\mathrm{fc}}
\newcommand{\BA}{B{\upharpoonright}A}
\newcommand{\BAc}{B{\upharpoonright}A^c}
\newcommand{\cc}{\mathfrak{c}}

% ---------- Title ----------
\title{CE Non-Derivability and the Ultralimit Representation Investigation:\\
\large Current Status}
\author{Patrick S.\ Mahon}
\date{2026-04-18}

\begin{document}
\maketitle

\begin{abstract}
This memorandum records the current state of two related investigations. The first
is complete: countable additivity is not first-order axiomatizable among finitely
additive probability Boolean algebras, and CE non-derivability follows as a
corollary. The second, spawned by that proof, asks which purely finitely additive
charges on a Boolean algebra arise as pointwise ultralimits of $\sigma$-additive
probabilities on the same algebra. That problem is settled in every tractable
regime — $\sigma$-algebras (negative), atomic algebras (positive), algebras embedding
join-preservingly into a standard measure algebra (positive), and algebras with too
few joins for $\sigma$-additivity to bite (positive vacuously) — and the remaining
open case is isolated as a single existence question, equivalent by a topological
reformulation to a problem about Radon-measure-free compact spaces. The pattern of
known results suggests this final question is independent of ZFC; further progress
will require forcing or consistency arguments.
\end{abstract}

\tableofcontents

\bigskip

\noindent\textit{Formal proof:} \texttt{ce\_nonderivability\_companion.tex}.\\
\textit{Archival working notes:} \texttt{finite\_cofinite\_calculation.md},
\texttt{rung3\_multidimensional\_pfa.md}, \texttt{stone\_geometric\_translation.md},
\texttt{row5\_candidate.md}, \texttt{subalgebra\_embedding\_lemma.md}.

\medskip\noindent\textit{Status tags used throughout:}
\textbf{Settled} (complete, written up),
\textbf{Eliminated} (approach shown not to work),
\textbf{Open} (genuinely unresolved).

% ==============================================================================
\section{Main Theorem}
% ==============================================================================

\begin{status}
\textbf{Settled.} Written up as a companion note targeting \textit{Annals of Pure
and Applied Logic}.
\end{status}

In the first-order language $\Lmu$ of Boolean algebras equipped with a normalized
finitely additive charge, no first-order theory characterizes those models whose
charge extends to a $\sigma$-additive measure on the generated $\sigma$-algebra.
Equivalently, countable additivity is not first-order axiomatizable among finitely
additive probability Boolean algebras.

The proof is a single ultraproduct argument. Dirac charges $\delta_n$ on the
finite-cofinite algebra $\mathcal{E}$ over $\Q$ are individually $\sigma$-additive.
Their ultraproduct along a non-principal ultrafilter $\mathcal{U}$ extending the
cofinite filter is the finite-cofinite charge $\ell$, which assigns mass~$1$ to
every cofinite set and mass~$0$ to every finite set. The decreasing sequence
$A_k = \Q\setminus\{q_0,\ldots,q_k\}$ has empty intersection yet $\ell(A_k)=1$
for all $k$, so $\ell$ is not $\sigma$-additive. By \L o\'s's theorem, any
first-order theory satisfied by all $\delta_n$ is satisfied by the ultraproduct;
but the ultraproduct fails $\sigma$-additivity. Contradiction.

\begin{corollary}[CE non-derivability]
No accumulation of first-order structural conditions on a directed system of Boolean
charge spaces implies collective exhaustion (CE). The finite-cofinite charge witnesses
failure of CE while satisfying every first-order condition.
\end{corollary}

\begin{remark}
The non-axiomatizability theorem relocates a foundational dispute from the domain of
terminological disagreement to the domain of logical necessity. The gap between
finitely additive coherence and genuine probability cannot be closed by any first-order
structural condition: some extra admissibility condition --- CE or an equivalent ---
is logically unavoidable.
\end{remark}

% ==============================================================================
\section{The Secondary Investigation}
% ==============================================================================

The proof above produces the finite-cofinite charge $\ell$ as an ultraproduct of
$\sigma$-additive charges on the same algebra $\mathcal{E}$. That raises a sharper
fixed-algebra question.

\begin{definition}[Ultralimit representability]
A finitely additive probability $\ell$ on a Boolean algebra $B$ is
\emph{ultralimit-representable} if there exists a sequence $(\mu_n)$ of
$\sigma$-additive probabilities on $B$ and a non-principal ultrafilter $\mathcal{U}$
on $\N$ such that
\[
  \ell(A) = \lim_{\mathcal{U}} \mu_n(A) \qquad \text{for all } A \in B.
\]
\end{definition}

The question is: for which Boolean algebras $B$ is every finitely additive
probability on $B$ ultralimit-representable?

This has a clean global dichotomy: on a $\sigma$-algebra the answer is never (the
Nikodym obstruction), on a general Boolean algebra it depends on the algebra's
structure. The investigation below classifies the answer by regime.

% ==============================================================================
\section{What Is Settled}
% ==============================================================================

% ------------------------------------------------------------------------------
\subsection{The \texorpdfstring{$\sigma$}{sigma}-algebra case}
% ------------------------------------------------------------------------------

\begin{status}
\textbf{Settled} (negative).
\end{status}

\begin{theorem}[Nikodym obstruction]
Let $\Sigma$ be a $\sigma$-algebra. No purely finitely additive charge on $\Sigma$
is ultralimit-representable.
\end{theorem}

\begin{proof}[Proof sketch]
Any pointwise ultralimit of $\sigma$-additive probabilities on a $\sigma$-algebra
is itself $\sigma$-additive, by the Nikodym--Vitali--Hahn--Saks convergence theorem
\cite{DunfordSchwartz1958}, IV.9.8. The non-trivial step is the interchange of the
ultralimit with the countable sum $\sum_k\mu_n(A_k)$; it does not follow from the
pointwise identity alone.
\end{proof}

\noindent
\textit{Consequence.} The representation question is entirely about Boolean algebras
that are not $\sigma$-complete. The finite-cofinite algebra is the minimal natural
setting where the Nikodym obstruction disappears.

% ------------------------------------------------------------------------------
\subsection{The atomic case}
% ------------------------------------------------------------------------------

\begin{status}
\textbf{Settled} (positive, uniformly).
\end{status}

If $B$ is atomic, every $\sigma$-additive probability on $B$ is a countable convex
combination of Dirac masses at atoms. Atoms of $B$ correspond to principal
ultrafilters $\Prin(B) \subseteq \St(B)$, and atomicity of $B$ is equivalent to
$\Prin(B)$ being dense in $\St(B)$. This density implies that every finitely
additive probability on $B$ is a pointwise limit of $\sigma$-additive ones.

\begin{theorem}[Atomic case]
If $B$ is atomic, every finitely additive probability on $B$ is ultralimit-representable.
\end{theorem}

This covers: the finite-cofinite algebra $\mathcal{E}$, all partition extensions
$\mathcal{E}_k$ (generated by $\mathcal{E}$ together with a finite $k$-partition of
$\N$ into infinite pieces), and $\mathcal{E}_\omega$ (countably many pieces). In
each case $\pfa(\mathcal{E}_k) \cong \Delta_{k-1}$, and the "spread mass uniformly
within each piece" construction gives the approximating sequence explicitly.

% ------------------------------------------------------------------------------
\subsection{The concrete non-atomic case}
% ------------------------------------------------------------------------------

\begin{status}
\textbf{Settled} (positive).
\end{status}

\begin{lemma}[Subalgebra embedding lemma]
Suppose $B$ embeds join-preservingly into a $\sigma$-algebra $\Sigma$ carrying a
strictly positive $\sigma$-additive probability. Then $\Supp(B) = \St(B)$: every
ultrafilter on $B$ lies in the support of some $\sigma$-additive probability on $B$.
In particular, every finitely additive probability on $B$ is ultralimit-representable.
\end{lemma}

Here \emph{join-preserving} means: whenever $\bigvee_n a_n$ exists in $B$, its
image in $\Sigma$ equals $\bigcup_n \iota(a_n)$. Restriction of the ambient
$\sigma$-additive probability to $B$ then directly inherits $\sigma$-additivity.

This covers all concrete non-$\sigma$-complete non-atomic algebras: interval
algebras, Borel subalgebras of standard measure spaces, and the clopen algebra
of any compact metrizable space carrying a strictly positive Borel measure.

\begin{corollary}
Every countably generated non-atomic Boolean algebra is ultralimit-representable.
\end{corollary}

\begin{proof}[Proof sketch]
A countably generated Stone space is compact metrizable. If $B$ is also non-atomic,
the Stone space has no isolated points. By Brouwer's theorem, a compact metrizable
zero-dimensional space with no isolated points is homeomorphic to the Cantor set.
The homeomorphism provides the join-preserving embedding into the Borel
$\sigma$-algebra of the Cantor set, and Lebesgue measure is strictly positive on
all clopens.
\end{proof}

% ------------------------------------------------------------------------------
\subsection{The vacuous case}
% ------------------------------------------------------------------------------

\begin{status}
\textbf{Settled} (positive, vacuously).
\end{status}

When $B$ has very few countable joins, $\sigma$-additivity becomes vacuous: the
condition $\mu(\bigvee_n a_n) = \sum_n\mu(a_n)$ only fires when $\bigvee_n a_n$
exists in $B$.

\begin{theorem}
Every finitely additive probability on $\mathcal{Q} = \mathcal{P}(\N)/\mathrm{fin}$
is $\sigma$-additive on $\mathcal{Q}$.
\end{theorem}

\begin{proof}[Proof sketch]
A pairwise disjoint sequence $(A_n)$ in $\mathcal{Q}$ with infinitely many nonzero
terms has no join: if $[S]$ were the join, then for each $n$ the selector element
$a_n \in A_n \cap S$ yields a strictly smaller upper bound $[S\setminus\{a_n\}]$,
contradicting minimality. The only disjoint sequences with an existing join have at
most finitely many nonzero terms, for which $\sigma$-additivity reduces to finite
additivity.
\end{proof}

% ------------------------------------------------------------------------------
\subsection{The canonical decomposition theorem}
% ------------------------------------------------------------------------------

\begin{status}
\textbf{Settled}.
\end{status}

\begin{theorem}[Canonical decomposition]
For any Boolean algebra $B$ and any element $A \in B$,
\[
  B \;\cong\; (\BA) \times (\BAc)
\]
canonically, where $\BA = \{C \in B : C \leq A\}$. The isomorphism sends
$C \mapsto (C \wedge A,\, C \wedge A^c)$.
\end{theorem}

If $\nu$ is $\sigma$-additive on $\BA$, then $\tilde\nu(C) = \nu(C \wedge A)$
is $\sigma$-additive on $B$: meets distribute over existing joins in any Boolean
algebra, so $(\bigvee_n C_n)\wedge A = \bigvee_n(C_n\wedge A)$ exists in $\BA$
whenever $\bigvee_n C_n$ exists in $B$, and $\sigma$-additivity transfers directly.

\begin{corollary}[Failures are direct-product failures]
If $\Supp(B) \subsetneq \St(B)$, then there exists a nonzero $A \in B$ such that
$\BA$ is measure-free and $B \cong (\BA)\times(\BAc)$.
\end{corollary}

\begin{proof}
Take any $A$ whose clopen $\hat A \subseteq \St(B)$ is disjoint from $\Supp(B)$.
Every $\sigma$-additive probability on $B$ assigns $\mu(A) = 0$. If $\BA$ admitted
a $\sigma$-additive probability $\nu$, the extension $\tilde\nu$ would satisfy
$\tilde\nu(A) = \nu(A) = 1$, contradicting $\mu(A)=0$ for all such $\mu$.
\end{proof}

\noindent
There is no such thing as an \emph{indecomposable} failure of full $\sigma$-additive
support: any failure is witnessed by a direct-product decomposition.

% ------------------------------------------------------------------------------
\subsection{The weak distributivity thread}
% ------------------------------------------------------------------------------

\begin{status}
\textbf{Settled} (does not apply to the non-$\sigma$-complete regime).
\end{status}

Dža\-monja--Plebanek~\cite{DzamonjaPleb2008} prove in ZFC: for a Boolean
\emph{$\sigma$-algebra}, weakly distributive plus strictly positive finitely additive
measure is equivalent to strictly positive $\sigma$-additive measure. The same
result appears in Fremlin~\cite{Fremlin2004}, \S391D, and
Plebanek~\cite{Plebanek2008}.

These theorems are sharp ZFC results in the $\sigma$-complete case. None applies
to non-$\sigma$-complete algebras: the proofs use countable suprema in essential
ways, and $\sigma$-additivity on a non-$\sigma$-complete algebra is already weaker
(it fires only when joins exist in $B$). Weak distributivity yields no necessary
condition on a Strategy~D counterexample (see Section~\ref{sec:frontier}).

% ==============================================================================
\section{What Was Ruled Out}
% ==============================================================================

\paragraph{Direct-product strategy via $\mathcal{Q}$.}
\begin{status}
\textbf{Eliminated}.
\end{status}
An earlier approach sought a negative example $B = \mathcal{I}\times\mathcal{Q}$
using a claimed fact that $\mathcal{Q} = \mathcal{P}(\N)/\mathrm{fin}$ supports no
$\sigma$-additive probability. That claim is false: every finitely additive
probability on $\mathcal{Q}$ is $\sigma$-additive (\S3.4). The strategy is
eliminated.

\paragraph{Indecomposable negative instances.}
\begin{status}
\textbf{Eliminated}.
\end{status}
An "indecomposable" negative instance would require $\Supp(B)\subsetneq\St(B)$
without any direct-product decomposition. The canonical decomposition theorem shows
this is impossible (\S3.5): any gap in $\sigma$-additive support forces a
direct-product decomposition with a measure-free factor.

\paragraph{Strategies A, B, C.}
\begin{status}
\textbf{Eliminated}.
\end{status}
Three construction strategies --- entangled Boolean extensions (A), quotient algebras
with measure-free tails (B), and forcing-generic constructions (C) --- all required
indecomposable negative instances. Since none can exist, all three are eliminated.

\paragraph{Weak distributivity as a non-$\sigma$-complete obstruction.}
\begin{status}
\textbf{Eliminated} (scope mismatch).
\end{status}
The Dža\-monja--Plebanek theorem was identified as a candidate source of a fifth
necessary condition. On careful reading it applies only to $\sigma$-complete
algebras; no analogous result is known for incomplete algebras. The thread is closed.

% ==============================================================================
\section{Current Frontier}
\label{sec:frontier}
% ==============================================================================

\paragraph{Strategy D.}
\begin{status}
\textbf{Open}.
\end{status}

\begin{quote}
\textit{Does there exist a non-$\sigma$-complete non-atomic Boolean algebra that
is measure-free (admits no $\sigma$-additive probability)?}
\end{quote}

If yes: pair it with any Boolean algebra carrying a $\sigma$-additive probability
to obtain a direct-product negative instance.

If no: any gap in $\sigma$-additive support would require a measure-free
non-$\sigma$-complete non-atomic factor, which would not exist; the open row
collapses and the representation problem is positive in all non-$\sigma$-complete
non-atomic cases.

\medskip

\noindent\textbf{Four necessary conditions on any counterexample.}
A Strategy~D counterexample must be: (a)~non-$\sigma$-complete, (b)~non-atomic,
(c)~measure-free, (d)~uncountably generated. Conditions (a)--(c) are by definition.
For~(d): every countably generated non-atomic Boolean algebra carries a strictly
positive $\sigma$-additive probability, by the Brouwer--Cantor reduction and the
subalgebra embedding lemma (\S3.3).

\medskip

\noindent\textbf{Topological reformulation.}
By Stone duality, Strategy~D is equivalent to:
\begin{quote}
\textit{Does there exist a compact totally disconnected Hausdorff space $K$ with no
isolated points, not basically disconnected, and carrying no strictly positive Radon
probability measure?}
\end{quote}
The dictionary is: no isolated points $=$ non-atomicity of $\mathrm{Clop}(K)$; not
basically disconnected $=$ $\mathrm{Clop}(K)$ is not a $\sigma$-algebra; no strictly
positive Radon measure $=$ measure-freeness of $\mathrm{Clop}(K)$. This connects
Strategy~D to the literature on Radon-measure-free compact spaces.

\medskip

\noindent\textbf{Set-theoretic sensitivity.}
The pattern of known results in that literature strongly suggests Strategy~D is
independent of ZFC.
\begin{itemize}[nosep]
\item Under $\mathrm{MA}+\neg\mathrm{CH}$: compact spaces of weight $<\cc$ carry
  strictly positive Radon measures~\cite{Fremlin2004}, \S531--534. Strategy~D
  counterexamples, if they exist, would need weight $\geq\cc$ under these axioms.
\item Under $\diamondsuit$: measure-free compact spaces of weight $\omega_1$ have
  been constructed \cite{Kunen1981,Fedorchuk1976}. Whether any such example is
  totally disconnected, non-atomic, and not basically disconnected remains to be
  checked.
\end{itemize}
No ZFC construction of a counterexample is known, and no ZFC proof that none exists.

% ==============================================================================
\section{What Would Move the Problem Next}
% ==============================================================================

\paragraph{Route 1 --- Forcing/consistency.}
Attempt to construct, under $\diamondsuit$ or a similar principle, a compact totally
disconnected non-atomic non-basically-disconnected measure-free space. Entry points:
Kunen~\cite{Kunen1981} and Fedorchuk~\cite{Fedorchuk1976} for measure-free compacta;
verify the non-atomic and non-basically-disconnected conditions for those
constructions. This is a set-theoretic topology project.

\paragraph{Route 2 --- Consistency of non-existence.}
Show, under $\mathrm{MA}+\neg\mathrm{CH}$, that every non-$\sigma$-complete
non-atomic Boolean algebra carries a strictly positive $\sigma$-additive probability.
Entry point: Fremlin~\cite{Fremlin2004}, \S531--534. This would establish the
positive half of an independence result.

\paragraph{Route 3 --- Targeted literature pass.}
The Radon-measure-free compact space literature may contain results that answer
Strategy~D directly, or reduce it further. The combined constraint of non-atomicity
and non-basic-disconnectedness alongside measure-freeness does not appear to have
been studied explicitly. A targeted pass through Fremlin vol.~5 and Plebanek's survey
articles on measures on Boolean algebras is cheaper than a forcing project and might
close the gap.

\paragraph{What is not useful.}
Further ZFC Boolean-algebra arguments. Weak distributivity. The
Dža\-monja--Plebanek theorem and its analogues. These threads are exhausted.

% ==============================================================================
\section{Summary}
% ==============================================================================

\begin{table}[h]
\centering
\begin{tabular}{@{}lll@{}}
\toprule
Regime & Answer & Method \\
\midrule
$\sigma$-complete non-atomic & \textbf{Negative} & Nikodym / Dunford--Schwartz IV.9.8 \\
Atomic & \textbf{Positive} & Density of $\Prin(B)$ in $\St(B)$ \\
Join-preserving into measure algebra & \textbf{Positive} & Subalgebra embedding lemma \\
Countably generated, non-atomic & \textbf{Positive} & Brouwer $+$ embedding lemma \\
Too few joins (e.g.\ $\mathcal{P}(\N)/\mathrm{fin}$) & \textbf{Positive (vacuous)} & $\sigma$-additivity $=$ finite additivity \\
Non-$\sigma$-complete non-atomic, intermediate joins & \textbf{Open} & Strategy~D \\
\bottomrule
\end{tabular}
\caption{Classification of the ultralimit representation problem by regime.}
\end{table}

The open row requires a measure-free direct-product factor that is itself
non-$\sigma$-complete and non-atomic. Strategy~D asks whether such a factor
can exist. The Boolean-algebraic methods employed throughout this investigation
do not reach that question; it lives at the set-theoretic and forcing level.

% ==============================================================================

\bibliographystyle{plain}
\bibliography{../../references}

\end{document}
